\chapter{Designing the Task REST API}
\chapterintro{Before writing SQL or handlers, define resources, representations, method semantics, validation, and errors. A written contract prevents implementation accidents from becoming the API.}

\begin{objectives}
\begin{itemize}
  \item Model a task resource and its collection representation.
  \item Define CRUD behavior with precise method and status semantics.
  \item Distinguish full replacement, partial modification, filtering, and pagination.
  \item Specify stable success and problem documents before implementation.
\end{itemize}
\end{objectives}

\section{Resource model}

The service exposes a task collection at \code{/tasks} and individual tasks at
\code{/tasks/\{id\}}. A task representation contains:

\begin{center}
\begin{tabularx}{\textwidth}{@{}>{\bfseries}p{27mm} p{29mm} X@{}}
\toprule
Field & Type & Rule \\
\midrule
id & integer & Server-generated positive identifier \\
title & string & Required, trimmed, 1--200 characters \\
description & string & Optional, at most 2000 characters; default empty \\
status & enum & \code{todo}, \code{in\_progress}, or \code{done} \\
priority & enum & \code{low}, \code{medium}, or \code{high} \\
due\_date & date or null & Exact \code{YYYY-MM-DD} form \\
created\_at & timestamp & Server-generated UTC RFC 3339 text \\
updated\_at & timestamp & Server-generated UTC RFC 3339 text \\
\bottomrule
\end{tabularx}
\end{center}

Clients may write only title, description, status, priority, and due date.
Unknown fields are rejected rather than ignored, catching misspellings and
preventing clients from believing a value was stored when it was not.

\section{Collection operations}

\method{POST /tasks} accepts a JSON object, applies defaults, creates a resource,
and returns \status{201 Created}. The body is the complete created representation
and \code{Location} identifies \code{/tasks/\{id\}}. The operation is not
idempotent: repeating it creates another task.

\method{GET /tasks} returns:

\begin{lstlisting}[style=wire]
{
  "items": [ ... task representations ... ],
  "total": 37,
  "limit": 20,
  "offset": 0
}
\end{lstlisting}

Supported filters are \code{status}, \code{priority}, \code{due\_before}, and
text query \code{q}. \code{limit} defaults to 20 and is between 1 and 100.
\code{offset} defaults to zero and is non-negative. Every supported parameter
may appear once; unknown or repeated parameters receive a validation problem.

Offset pagination is simple and suitable for a teaching database. Under
concurrent inserts it can repeat or skip items between pages. Cursor pagination
is often preferable for large, changing datasets but requires a stable ordering
key and a documented cursor format.

\section{Individual operations}

\method{GET /tasks/\{id\}} returns 200 with a task or 404 with a stable
\code{task-not-found} problem.

\method{PUT /tasks/\{id\}} is full replacement. Title is required. Omitted
optional fields return to defaults. Server-owned identity and creation timestamp
remain unchanged; update timestamp changes. Repeating the same PUT has the same
intended resource effect, although the timestamp changes.

\method{PATCH /tasks/\{id\}} updates only supplied writable fields. An empty
object is rejected. PATCH is not assumed idempotent in general; this particular
merge-style field assignment is idempotent for the same values.

\method{DELETE /tasks/\{id\}} deletes an existing task and returns 204 with no
content. Deleting an absent task returns 404. The intended final resource state
still makes DELETE idempotent.

\section{Validation sequence}

Validation proceeds from transport toward domain:

\begin{enumerate}
  \item HTTP parser validates message framing and limits.
  \item A body-required endpoint requires \code{Content-Length} and nonempty bytes.
  \item Request JSON helper requires the media type, UTF-8, and JSON grammar.
  \item Task validation requires an object, writable names, types, lengths,
        enum membership, and exact date format.
  \item The repository uses database constraints as a final durable backstop.
\end{enumerate}

Validation errors use 422 and an \code{errors} object. Malformed JSON remains
400 because no valid JSON value exists to validate. Unsupported media type is
415. These distinctions help a client decide what to fix.

\begin{wireview}
\begin{wiremessage}
POST /tasks HTTP/1.1
Host: 127.0.0.1:8080
Content-Type: application/json
Content-Length: 47

{"title":"Learn routing","priority":"high"}
\end{wiremessage}
The server supplies \code{status=todo}, an empty description, null due date,
timestamps, and an identifier.
\end{wireview}

\section{Error contract}

Every application error is \code{application/problem+json}. A task validation
problem uses a stable type URI. A client can render the title/detail for a person
and map \code{errors} to form fields. It should use \code{status} and \code{type}
for program logic.

Not every failure is safe to expose. A database exception becomes a generic 500
problem; the traceback is logged. A missing task is ordinary control flow and
does not produce an error log.

\section{What the API deliberately omits}

There are no users, projects, labels, attachments, authentication, authorization,
soft deletion, synchronization, ETags, idempotency keys, or schema negotiation.
Each could be valuable, but each introduces new resource and security decisions.
The capstone stays focused on HTTP, framework flow, validation, SQL, and tests.

\summarytitle

The contract defines tasks, writable fields, collection and individual methods,
filter/pagination rules, validation layers, and structured failures. PUT replaces,
PATCH partially updates, POST creates, and DELETE removes. Implementation can now
be judged against a stable design.

\reviewtitle
\begin{enumerate}
  \item Why does POST return both a representation and Location?
  \item What happens to omitted optional fields in PUT versus PATCH?
  \item Why can offset pagination become inconsistent under concurrent inserts?
  \item Which layer distinguishes malformed JSON from invalid task fields?
\end{enumerate}

\begin{exercise}[Contract examples]
Write one complete request and response example for every endpoint, including
headers and status. Add examples for every error class. These become documentation
fixtures and integration-test candidates.
\end{exercise}

\begin{exercise}[Extend without coding]
Design an optional \code{completed\_at} field. Define who writes it, how status
transitions affect it, PUT/PATCH behavior, database constraints, filtering, and
backward compatibility. Notice how one field crosses every layer.
\end{exercise}

